\documentclass{article}

% Packages required to support encoding
\usepackage{ucs}
\usepackage[utf8x]{inputenc}

% Packages required by code


% Packages always used
\usepackage{hyperref}
\usepackage{xspace}
\usepackage[usenames,dvipsnames]{color}
\hypersetup{colorlinks=true,urlcolor=blue}



\begin{document}

\vspace{.5em} \hrule \vspace{.5em}
title: Python Tricks For Data Scientist layout: single mathjax: true category: Data-Science

\hypertarget{tags_data_science_python}{}\subsection*{{tags: data science, python}}\label{tags_data_science_python}

Recently, I've found myself gravitating toward a small set of `tricks' in Python over and over again and thought it would be fun to share and discuss a few of them with other data scientists. So here they go:

\hypertarget{cached_data_loader}{}\subsection*{{Cached Data Loader}}\label{cached_data_loader}

Jupyter notebooks are great for exploratory data analysis for many reasons. One major of which is their persistent state between code execution: You can load your data once and then iteratively perform all sorts of modifications / visualizations without having to reload your data over the network.

Sometimes, however, that's still not enough: If your network connection is slow or the dataset you're trying to analyze particular large or you need to restart your kernel frequently for whatever reason, data loading can become a real pain.

To this, I've found myself using cached data loading more and more frequently in recent times. Here's how this may look like in practice:

``thon def \emph{cached}data\_loader( connector: StorageConnector, query: str, query\_parameters: ListQueryParameter, cache\_location: Path, force\_update: bool, ) -{\tt \symbol{62}} pandas.DataFrame: if force\_update or not os.path.isfile(cache\_location): data = connector.query\_to\_dataframe( query, query\_parameters ) data.to\_pickle(cache\_location) else: data = pd.from\_pickle(cache\_location)

\begin{verbatim}return data\end{verbatim}
def load\_marketing\_data( query\_parameters: ListQueryParameter = , cache\_location: Path = Path(``cache/marketing.pkl''), force\_update: bool = False ) -{\tt \symbol{62}} pandas.DataFrame: connector = \ldots{} query = ``\ldots{}'' return = \emph{cached}data\_loader( connector, query, query\_parameters, cache\_location, force\_update )

``

After downloading marketing data over the network, any subsequent call to this function will load a cached version from disk instead. However, by setting {\colorbox[rgb]{1.00,0.93,1.00}{\tt force\char95update\char32\char61\char32True}} you can still refresh your (cached) data at any time.

This can probably be implemented much more elegantly as a decorator. And if you do, I'd love to hear about it!

\hypertarget{debugging_with_upturn}{}\subsection*{{Debugging with Upturn}}\label{debugging_with_upturn}

Since many bugs in data science code only occur on certain input data, debugging can be challenging and simple debugging tools often seem too restrictive, whereas more feature complete debugging tools can be intimidating. What I've found to be useful at times is to abuse IPython as a debugger. Say, for instance, that a certain pre-processing step in my data pipeline yields unexpected results. Instead of setting a breakpoint and attempting to understand what's going on with a debugger, I might add something like this:

``thon def preprocessing(raw\_data: pd.DataFrame) -{\tt \symbol{62}} np.ndarray: \ldots{} buggy code \ldots{}

\begin{verbatim}import IPython
IPython.embed()
exit(1)

...
return data\end{verbatim}
``

When reaching {\colorbox[rgb]{1.00,0.93,1.00}{\tt IPython\char46embed\char40\char41}}, this opens an interactive Python session in the terminal that allows me to inspect all relevant variables, execute arbitrary code, experiment with possible solutions, \ldots{} When you're done, simply exit IPython which will then also exit the program in the following line.

\hypertarget{diffing_and_merging_jupyter_notebooks}{}\subsection*{{Diffing and Merging Jupyter Notebooks}}\label{diffing_and_merging_jupyter_notebooks}

Working in close collaboration on challenging problems is easily my favorite aspect of being a data scientist. Resolving merge conflicts in Jupyter notebooks, however, that close collaboration inevitably leads to, has been a rather dreadful experience for me. That is, until I found \href{https://github.com/jupyter/nbdime}{nbdime}. It's a super handy tool to diff and merge Jupyter notebooks in your terminal or via their web-based GUI.



\hypertarget{speeding_up_bigquery}{}\subsection*{{Speeding Up BigQuery}}\label{speeding_up_bigquery}

Transforming large amounts of data from \href{https://cloud.google.com/bigquery}{BigQuery} to Pandas DataFrames, even from within Google's Compute Engine, can be incredibly slow. Fortunately, Google offers an alternative that I've personally experienced to shorten load times by two orders of magnitude -- BigQuery's Storage API. Here's how to use it (assuming that this API has been enabled in your project and your service account has permissions to use it):

``thon from google.cloud import bigquery, bigquery\_storagev1beta1

credentials, project\_id = \ldots{}

bq\_client = bigquery.Client( credentials=credentials, project=project\_id )

bq\_storage\_client = bigquery\_storagev1beta1.BigQueryStorageClient( credentials=credentials )

query = ``\ldots{}''

data = ( bq\_client.query(query) .result() .to\_dataframe(bq\_storage\_client=bq\_storage\_client) )

``

Keep in mind, however, that this will result in additional costs.


\end{document}
